\documentclass[10pt,a4paper]{amsart}
\usepackage[margin=19mm]{geometry}
\usepackage{amsmath,amssymb,mathtools}
\usepackage[scr=rsfs,cal=euler]{mathalfa}
\usepackage{xfrac}
\usepackage[unicode,hidelinks,bookmarks=false]{hyperref}
\newcommand{\ie}{i.e.\ }
\newcommand{\cf}{cf.\ }
\newcommand{\resp}{respectively\ }
\newcommand{\Subsection}{\subsection}
\providecommand{\nobreakdash}{\nobreak\hbox{-}}
\setlength{\emergencystretch}{2em}
\multlinegap0pt
\usepackage[shortlabels]{enumitem}
\usepackage{dsfont}
\usepackage{amssymb}
\usepackage{algorithm}\usepackage{algpseudocode}
\usepackage{mathtools}
\usepackage{tikz}
\usepackage{bm}
\usepackage{tcolorbox}
\newtheorem{assumption}{Assumption}
\newtheorem{proposition}{Proposition}
\newtheorem{lemma}{Lemma}
\newcommand{\D}[1]{\mbox{\rm #1}} 
\newcommand{\dd}{\D{d}}
\title{A derivative-free particle method for optimization in Hilbert spaces}
\author{Hui Huang, Hicham Kouhkouh}
\date{Source: arXiv:2605.31565v1 (CC BY 4.0)}
\begin{document}
\maketitle

\noindent\emph{Source and licence.} A derivative-free particle method for optimization in Hilbert spaces. Version arXiv:2605.31565v1 (CC BY 4.0), distributed by arXiv under the Creative Commons Attribution 4.0 International licence, http://creativecommons.org/licenses/by/4.0/, as recorded in the arXiv licence metadata for this exact version and shown on the abstract page https://arxiv.org/abs/2605.31565v1. Full licence notice: \texttt{licenses/arxiv-2605.31565-CC-BY-4.0.txt}.\par\smallskip

\noindent\textit{Editorial scope.} Counted unit: the article's lemma F (source0.tex lines 1654--1656). The article's complete original proof of it is delivered with it (source0.tex lines 1658--1810, one \texttt{proof} environment of 153 lines). No mathematics is written, repaired or re-derived: the statement and the proof are the article's own lines, in the article's own notation, and no retained line is substituted or re-flowed.  Retained uncounted as the source-local context the proof needs: lines 466--491; lines 495--502; lines 586--599; lines 603--608; lines 619--645; lines 1603--1610; lines 1612--1614.  Nothing was omitted from any retained span, so the package carries no omission record.  No retained line contains a citation, so the wrapper carries no bibliography and no bibliography entry.  No retained line contains a figure environment, an image-inclusion command or drawing code, so no image is delivered and the package declares no asset dependency.  Editorial formatting only, in addition: an AMS \texttt{amsart} preamble that restates the article's own declarations and macros used by the retained lines, namely the environments \texttt{assumption}, \texttt{proposition}, \texttt{lemma} on separate counters, together with the standard packages those lines need.  The retained environments print in this document as Assumption 1, Proposition 1, Lemma 1 in order of appearance, whereas the article prints the same items with the numbering of its own full document.  The complete native source of the article as distributed in the arXiv e-print for this exact version is bundled unchanged as \texttt{sources/201732/source0.tex}.
\begin{assumption}[Regularity of the Objective Function]
\label{assum:objective}
The function $\mathscr{E}$ satisfies the following properties:
\begin{enumerate}
    \item \textbf{Bounded from below:} Without loss of generality, we assume 
    \begin{equation*}
        \inf_{x \in H} \mathscr{E}(x) \geq 0 .
    \end{equation*}
    If bounded below by a negative constant, then a simple shift would leave the consensus point $\mathfrak{m}_\alpha$ invariant. 
    \item \textbf{Polynomial growth:} There exist constants $C_{\mathscr{E}} > 0$ and $p \geq 1$ such that 
    \begin{equation*}
        \mathscr{E}(x) \leq C_{\mathscr{E}} \big(1 + \|x\|_H^p\big) \qquad \forall\, x\in H.
    \end{equation*}
    \item  \textbf{Coercive with sufficient growth:} There exist constants $c, R_0 > 0$ and $p \geq 1$ such that 
       \begin{equation*}
       \mathscr{E}(x) \geq c\|x\|_H^p \quad \mbox{ for all }\quad \|x\|_H > R_0.
    \end{equation*}
    \item \textbf{Local Lipschitz continuity:} For any radius $R > 0$, there exists a constant $L > 0$ such that for all $x, y \in H$ and $p\geq 1$,
    \begin{equation*}
    \begin{aligned}
        |\mathscr{E}(x) - \mathscr{E}(y)|
        & \leq L(1+\|x\|_H^{p-1}+\|y\|_H^{p-1})\|x - y\|_H.
    \end{aligned}
    \end{equation*}
\end{enumerate}
\end{assumption}

\begin{proposition}
\label{prop:global_lipschitz}
Under Assumption \ref{assum:objective}, the mappings 
\begin{equation*}
    w(x) = \exp(-\alpha\, \mathscr{E}(x)) \in \mathbb{R}, \qquad \text{ and } \qquad f(x) = x\,\exp(-\alpha\, \mathscr{E}(x)) \,\in H
\end{equation*}
are globally Lipschitz continuous on $H$.
\end{proposition}

\begin{assumption}[Finite-Dimensional Active Subspace] \label{assump:active_subspace}
We assume there exists a linear subspace $V \subset H$ satisfying the following conditions:

\begin{enumerate}[label=(\roman*)]
    \item 
    The subspace $V$ has a finite dimension $d < \infty$, and the ambient space admits the orthogonal decomposition $H = V \oplus V^\perp$. Furthermore, the global minimizer $x^*$ of $\mathscr{E}$ resides in $V$, and the initial probability measure $\mu_0 = \mathrm{Law}(\xi)$ is supported entirely on $V$. 
    
    \item 
    The restriction of $Q$ to the active subspace $V$ is strictly positive-definite, that is there exists a constant $\lambda_{\min} > 0$ such that
    \begin{equation*}
        \langle v, Q v \rangle_H \geq \lambda_{\min} \|v\|_H^2 \qquad \forall\,v\in V.
    \end{equation*}
\end{enumerate}
\end{assumption}

\begin{equation}\label{def:diffusion}
\begin{aligned}
    \mathsf{D}: H & \longrightarrow \mathscr{L}(H,H) \\
    x & \longmapsto \mathsf{D}(x) := \|x\|_H \mathsf{P}_V\,,
\end{aligned}
\end{equation}

\begin{proposition}
\label{prop:diffusion}
Let Assumption \ref{assump:active_subspace} hold, and $V\subset H$ be as defined therein.  
Then the mapping $\mathsf{D}: H \to \mathscr{L}_Q(H, H)$ defined in \eqref{def:diffusion} satisfies the following properties:

\begin{enumerate}
    \item 
    For all $x, y \in H$, it holds that
    \begin{equation*}
        \|\mathsf{D}(x) - \mathsf{D}(y)\|_{\mathscr{L}_Q} \leq   L_{\mathsf{D}} \|x - y\|_H,
    \end{equation*}
    where $L_{\mathsf{D}}:=\|\mathsf{P}_V\|_{\mathscr{L}_Q}=\mathrm{Tr}(\mathsf{P}_V Q \mathsf{P}_V)^{1/2}$.
    
    \item 
    For all $x\in H$
    \begin{equation*}
        \|\mathsf{D}(x)\|_{\mathscr{L}_Q} \leq  L_{\mathsf{D}}\|x\|_H \qquad \text{ and } \qquad \mathsf{D}(0) = 0.
    \end{equation*}

    
    \item  
    For all $x \in H$ and all $v \in V$
    \begin{equation*}
        \|\mathsf{D}(x)^* v\|_H^2 =  \|x\|_H^2 \|v\|_H^2.
    \end{equation*}
\end{enumerate}
\end{proposition}

\begin{equation}
\label{eq:mathds F}
\begin{aligned}
    X^{i} - \mathtt{G}(\mathds{X}) 
    & =  X^{i} - \frac{1}{\sum\limits_{k=1}^{N}\omega_{\alpha}(X^{k})} \, \sum\limits_{j=1}^{N}\,X^{j}\; \omega_{\alpha}(X^{j})\\ 
    & = \sum\limits_{j=1}^{N}\,\frac{\big(X^{i}-X^{j}\big)\; \omega_{\alpha}(X^{j})}{\sum\limits_{k=1}^{N}\omega_{\alpha}(X^{k})} \quad =: \; \mathbf{F}^{i}(\mathds{X}) \in \mathbb{H}.
\end{aligned}
\end{equation}

\begin{equation}\label{CBO in H N F}
    \dd \mathds{X}_{t} = -\lambda \, \mathds{F}(\mathds{X}_{t}) \, \dd t + \sigma\, \mathbb{D}(\mathds{F}(\mathds{X}_{t})) \, \dd \mathbb{W}^{Q}_{t},\quad t>0, \quad \mathds{X}_{0}=\mathbf{x}_{\circ}=(\mathbf{x}^{i}_{\circ})_{i\in[N]},
\end{equation}

\begin{lemma}\label{lem: F}
Let Assumption \ref{assum:objective} and Assumption \ref{assump:active_subspace} be satisfied, and suppose we are in the setting of Proposition \ref{prop:diffusion}. Then the drift and diffusion coefficients of \eqref{CBO in H N F} are locally Lipschitz continuous. In particular, the map $\mathds{F}:H^N \to H^N$ is locally Lipschitz continuous.
\end{lemma}

\begin{proof}
Thanks to item (1) of Proposition \ref{prop:diffusion}, $\mathsf{D}$ is Lipschitz, hence also $\mathbb{D}$. It suffices then to show that $\mathds{F}$ is Lipschitz. 
Let us recall 
\begin{equation*}
\begin{aligned}
    \mathds{F}: H^N & \longrightarrow H^N,\;
        \mathds{X}  \longmapsto \big( \mathbf{F}^{i}(\mathds{X}) \big)_{i\in [N]} %\in H^{N}
\end{aligned}
\end{equation*}
and for each $i\in [N]$, we have \eqref{eq:mathds F} that is
\begin{equation*}
\begin{aligned}
    \mathbf{F}^{i}: H^N & \longrightarrow H,\quad
        \mathds{X}  \longmapsto \sum\limits_{j=1}^{N}\,\frac{\big(X^{i}-X^{j}\big)\; \omega_{\alpha}(X^{j})}{\sum\limits_{k=1}^{N}\omega_{\alpha}(X^{k})}
\end{aligned}
\end{equation*}
It is therefore sufficient to verify that $\mathbf{F}^i$ is globally Lipschitz continuous. 



Let $i\in[N]$ be arbitrarily fixed. 
Let $\mathds{X}, \mathds{Y} \in H^N$ satisfy $\;\|\mathds{X}\|_{H^N}, \|\mathds{Y}\|_{H^N} \le R\;$
for some fixed $R>0$. 
We shall estimate the quantity 
\begin{equation*}
    \mathbf{F}^{i}(\mathds{X}) - \mathbf{F}^{i}(\mathds{Y}) \, \in \mathbb{H},\quad \forall\, i\in [N].
\end{equation*}
We define
\begin{equation*}
    A^i(\mathds{X}) := \sum_{j=1}^N (X^i - X^j)\,\omega_\alpha(X^j), \qquad B(\mathds{X}) := \sum_{k=1}^N \omega_\alpha(X^k), \qquad \text{ so that } \; \mathbf{F}^i(\mathds{X}) = \frac{A^i(\mathds{X})}{B(\mathds{X})}.
\end{equation*}
Let us also introduce
\begin{equation*}
    M_\omega(R) := \sup_{\|x\|_H \le R} \omega_\alpha(x), \qquad m_\omega(R) := \inf_{\|x\|_H \le R} \omega_\alpha(x) > 0.
\end{equation*}
and define
\begin{equation*}
    M_B(R) := N M_\omega(R), \qquad m_B(R) := N m_\omega(R).
\end{equation*}

\textbf{Step 1.} \textit{(The quantities depending on $B$)}  
For all $\mathds{X}$ such that $\|\mathds{X}\|_{H^N} \le R$, it holds
\begin{equation}\label{eq: B bound}
m_B(R) \le B(\mathds{X}) \le M_B(R).
\end{equation}
On the other hand, recalling Proposition \ref{prop:global_lipschitz}, we have $\omega_\alpha$ is globally Lipschitz with a constant denoted $L_\omega$. Therefore
\begin{equation*}
    |B(\mathds{X}) - B(\mathds{Y})| \le \sum_{k=1}^N |\omega_\alpha(X^k) - \omega_\alpha(Y^k)| \le L_\omega \sum_{k=1}^N \|X^k - Y^k\|_H \le L_\omega \|\mathds{X} - \mathds{Y}\|_{H^N}.
\end{equation*}
Moreover, 
\begin{equation}\label{eq: B lip}
    \left|\frac{1}{B(\mathds{X})} - \frac{1}{B(\mathds{Y})}\right| = \frac{|B(\mathds{X}) - B(\mathds{Y})|}{B(\mathds{X}) B(\mathds{Y})} \le \frac{L_\omega}{m_B(R)^2} \|\mathds{X} - \mathds{Y}\|_{H^N}.
\end{equation}

\textbf{Step 2.} \textit{(The quantities depending on $A$)} 
We estimate 
\begin{equation*}
    A^i(\mathds{X}) - A^i(\mathds{Y}) = \sum_{j=1}^N \Big[(X^i - X^j)\omega_\alpha(X^j) - (Y^i - Y^j)\omega_\alpha(Y^j)\Big].
\end{equation*}
We split each term as
\begin{equation*}
\begin{aligned}
    &(X^i - X^j)\omega_\alpha(X^j) - (Y^i - Y^j)\omega_\alpha(Y^j) \\
    & = (X^i - X^j - Y^i + Y^j)\omega_\alpha(X^j) + (Y^i - Y^j)(\omega_\alpha(X^j) - \omega_\alpha(Y^j)).
\end{aligned}
\end{equation*}
\textbullet\; Since $\omega_\alpha$ is bounded on bounded sets, we have
\begin{equation*}
    \|(X^i - X^j - Y^i + Y^j)\omega_\alpha(X^j)\|_H \le M_\omega(R)\big(\|X^i - Y^i\|_H + \|X^j - Y^j\|_H\big).
\end{equation*}
\textbullet\; Since $\|\mathds{Y}\|_{H^N} \le R$, we have
\begin{equation*}
    \|Y^i - Y^j\|_H \le 2R.
\end{equation*}
\textbullet\; Using the Lipschitz property of $\omega_\alpha$ with constant $L_\omega$,
\begin{equation*}
    |\omega_\alpha(X^j) - \omega_\alpha(Y^j)| \le L_\omega \|X^j - Y^j\|_H,
\end{equation*}
thus we obtain
\begin{equation*} 
    \|(Y^i - Y^j)(\omega_\alpha(X^j) -    \omega_\alpha(Y^j))\|_H \le 2R L_\omega \|X^j - Y^j\|_H.
\end{equation*}

Therefore, for each $j$,
\begin{equation*}
\begin{aligned}
    &\|(X^i - X^j)\omega_\alpha(X^j) - (Y^i - Y^j)\omega_\alpha(Y^j)\|_H \\
    & \le M_\omega(R)\|X^i - Y^i\|_H + \big(M_\omega(R) + 2R L_\omega\big)\|X^j - Y^j\|_H.
\end{aligned}
\end{equation*}
Summing over $j=1,\dots,N$ yields
\begin{equation*}
    \|A^i(\mathds{X}) - A^i(\mathds{Y})\|_H \,\le\, N\, M_\omega(R)\|X^i - Y^i\|_H + N\big(M_\omega(R) + 2R L_\omega\big) \|\mathds{X} - \mathds{Y}\|_{H^N}.
\end{equation*}
Finally, using $\|X^i - Y^i\|_H \le \|\mathds{X} - \mathds{Y}\|_{H^N}$, we obtain
\begin{equation*}
    \|A^i(\mathds{X}) - A^i(\mathds{Y})\|_H \le N\big(2R M_\omega(R) + 2R L_\omega\big) \|\mathds{X} - \mathds{Y}\|_{H^N}.
\end{equation*}
Let us define the constant
\begin{equation*}
    C_A(R) := N\big(2R M_\omega(R) + 2R L_\omega\big).
\end{equation*}
Thus,
\begin{equation}\label{eq: A lip}
    \|A^i(\mathds{X}) - A^i(\mathds{Y})\|_{H^N} \le C_A(R)\|\mathds{X} - \mathds{Y}\|_{H^N}.
\end{equation}

On the other hand, let us recall that
\begin{equation*}
    A^i(\mathds{Y}) = \sum_{j=1}^N (Y^i - Y^j)\,\omega_\alpha(Y^j).
\end{equation*}
Taking the $H$-norm and using the triangle inequality gives
\begin{equation*}
    \|A^i(\mathds{Y})\|_{H} \le \sum_{j=1}^N \|Y^i - Y^j\|_{H}\,\omega_\alpha(Y^j).
\end{equation*}
Since $\|\mathds{Y}\|_{H^N} \le R$, we have $\;\|Y^i - Y^j\|_{H} \le \|Y^i\|_{H} + \|Y^j\|_{H} \le 2R.\;$
Moreover, $\omega_\alpha$ is bounded on bounded sets, so  $\;\omega_\alpha(Y^j) \le M_\omega(R).\;$ Combining these estimates yields
\begin{equation*}
    \|Y^i - Y^j\|_{H}\,\omega_\alpha(Y^j) \le 2R\,M_\omega(R).
\end{equation*}
Summing over $j=1,\dots,N$, we obtain
\begin{equation}\label{eq: A bound}
    \|A^i(\mathds{Y})\|_{H} \le \sum_{j=1}^N 2R\,M_\omega(R) = N \, 2R \, M_\omega(R).
\end{equation}

\textbf{Step 3.} \textit{(Conclusion)} 
We write
\begin{equation*}
    \mathbf{F}^i(\mathds{X}) - \mathbf{F}^i(\mathds{Y}) = \frac{A^i(\mathds{X}) - A^i(\mathds{Y})}{B(\mathds{X})} + A^i(\mathds{Y}) \left( \frac{1}{B(\mathds{X})} - \frac{1}{B(\mathds{Y})} \right).
\end{equation*}
\textbullet\; Using previous bounds \eqref{eq: B bound} and \eqref{eq: A lip},
\begin{equation*}
    \left\|\frac{A^i(\mathds{X}) - A^i(\mathds{Y})}{B(\mathds{X})}\right\|_{H^N} \le \frac{C_A(R)}{m_B(R)} \|\mathds{X} - \mathds{Y}\|_{H^N}.
\end{equation*}
\textbullet\; Using previous bounds \eqref{eq: B lip} and \eqref{eq: A bound},
\begin{equation*}
    \|A^i(\mathds{Y})\| \left|\frac{1}{B(\mathds{X})} - \frac{1}{B(\mathds{Y})}\right| \le \frac{2N R M_\omega(R)\, L_\omega}{m_B(R)^2} \|\mathds{X} - \mathds{Y}\|_{H^N}.
\end{equation*}
Thus we obtain
\begin{equation*}
    \|\mathbf{F}(\mathds{X}) - \mathbf{F}(\mathds{Y})\|_{H^N} \le C(R)\|\mathds{X} - \mathds{Y}\|_{H^N},
\end{equation*}
where $C(R)$ is a constant defined by
\begin{equation*}
    C(R) := \frac{C_A(R)}{m_B(R)} + \frac{2N R M_\omega(R)\, L_\omega}{m_B(R)^2} =  \frac{2R M_\omega(R) + 2R L_\omega}{m_\omega(R)} + \frac{2 R M_\omega(R)\, L_\omega}{N(m_\omega(R))^2}
\end{equation*}
Hence $\mathbf{F}$ is locally Lipschitz on $H^N$.

To show the linear growth of $\mathbf{F}^{i}$, it suffices to observe that
\begin{equation*}
    |\mathbf{F}^{i}(\mathds{X})| \leq \left|  X^{i} - \frac{1}{\sum\limits_{k=1}^{N}\omega_{\alpha}(X^{k})} \, \sum\limits_{j=1}^{N}\,X^{j}\; \omega_{\alpha}(X^{j})  \right| \leq |X^{i}| + \|\mathds{X}\|\leq 2\,\|\mathds{X}\|.
\end{equation*}
\end{proof}

\end{document}
